%============================================================
%  preamble.tex  –  packages, settings, custom commands
%============================================================

% ---- Encoding & language ----------------------------------
\usepackage[T1]{fontenc}
\usepackage[utf8]{inputenc}
\usepackage[english]{babel}

% ---- Fonts ------------------------------------------------
\usepackage{lmodern}          % Better Computer Modern
\usepackage{microtype}        % Microtypographic improvements

% ---- Page geometry ----------------------------------------
\usepackage[
  a4paper,
  top=25mm, bottom=30mm,
  inner=30mm, outer=25mm,
  headheight=15pt
]{geometry}

% ---- Headers & footers ------------------------------------
\usepackage{fancyhdr}
\pagestyle{fancy}
\fancyhf{}
\fancyhead[LE]{\small\leftmark}
\fancyhead[RO]{\small\rightmark}
\fancyfoot[C]{\thepage}
\renewcommand{\headrulewidth}{0.4pt}

% ---- Mathematics ------------------------------------------
\usepackage{amsmath, amssymb, amsthm}
\usepackage{mathtools}        % Extends amsmath
\usepackage{bm}               % Bold math \bm{}
\usepackage{physics}          % Dirac notation \bra{}, \ket{}, \braket{}
\usepackage{siunitx}          % SI units

% ---- Code snippets ----------------------------------------
\newcommand{\code}[1]{\lstinline[language=Python]|#1|}

% ---- Chemistry --------------------------------------------
\usepackage[version=4]{mhchem}  % Chemical formulae \ce{H2O}
\usepackage{chemfig}             % Structural formulae (optional)

% ---- Graphics & colors ------------------------------------
\usepackage{graphicx}
\usepackage[dvipsnames, svgnames, table]{xcolor}

% KU brand colour  (Pantone 485 C ≈ #901A1E)
\definecolor{KURed}{HTML}{901A1E}
\definecolor{KUDarkBlue}{HTML}{003366}

\graphicspath{{figures/}}

% ---- Figures & tables -------------------------------------
\usepackage{booktabs}         % Professional tables
\usepackage{multirow}
\usepackage{tabularx}
\usepackage{longtable}
\usepackage{caption}
\usepackage{subcaption}
\usepackage{float}

% ---- Bibliography -----------------------------------------
\usepackage[numbers, sort&compress]{natbib}

% ---- Hyperlinks -------------------------------------------
\usepackage[
  colorlinks = true,
  linkcolor  = KURed,
  citecolor  = KURed,
  urlcolor   = KUDarkBlue
]{hyperref}
% PDF metadata set via \hypersetup to avoid keyval comma-parsing issues
\hypersetup{
  pdfauthor  = {Tobias},
  pdftitle   = {Optimized Virtual Orbital Space for Quantum Computing},
  pdfsubject = {Master Thesis -- University of Copenhagen}
}
\usepackage{cleveref}         % \cref{} – auto-labels (load after hyperref)

% ---- Code listings ----------------------------------------
\usepackage{listings}
\usepackage{algorithm}
\usepackage{algpseudocode}

\lstdefinestyle{python}{
  language     = Python,
  basicstyle   = \ttfamily\small,
  keywordstyle = \color{KURed}\bfseries,
  commentstyle = \color{gray}\itshape,
  stringstyle  = \color{teal},
  numberstyle  = \tiny\color{gray},
  numbers      = left,
  frame        = single,
  breaklines   = true,
  showstringspaces = false,
}

% ---- Acronyms --------------------------------------------
\usepackage{acronym}          % \ac{}, \acro{}, \acrodef{}

% ---- Miscellaneous ----------------------------------------
\usepackage{enumitem}         % Control list spacing
\usepackage{csquotes}         % Context-sensitive quotation marks
\usepackage{lipsum}           % Dummy text (remove in final version)
\usepackage{appendix}
\usepackage{emptypage}        % suppress headers on blank pages
\usepackage{afterpage}

% ---- Custom theorem environments --------------------------
\theoremstyle{definition}
\newtheorem{definition}{Definition}[chapter]
\newtheorem{theorem}{Theorem}[chapter]
\newtheorem{proposition}{Proposition}[chapter]

% ---- Custom commands: quantum chemistry shortcuts ---------
\newcommand{\Hhat}{\hat{H}}
\newcommand{\Etot}{E_{\mathrm{tot}}}
\newcommand{\Ecorr}{E_{\mathrm{corr}}}
\newcommand{\EHF}{E_{\mathrm{HF}}}
\newcommand{\EMP}[1]{E^{(\mathrm{MP#1})}}
\newcommand{\Jtwo}{J_2}
\newcommand{\ovos}{\textsc{ovos}}
\newcommand{\pyscf}{\textsc{PySCF}}
\newcommand{\uhf}{\textsc{UHF}}
\newcommand{\rhf}{\textsc{RHF}}
\newcommand{\MPn}[1]{\textsc{MP#1}}  % use \MPn{2} for MP2 etc.; \mp is reserved (∓)
\newcommand{\ccsd}{\textsc{CCSD}}

% Two-electron integral in physicist's notation
\newcommand{\braketij}[2]{\langle #1 \| #2 \rangle}   % antisymmetrised
\newcommand{\tamp}[4]{t_{#1 #2}^{#3 #4}}               % MP1 amplitude

% ---- Chapter/section styling ------------------------------
\usepackage{titlesec}
\titleformat{\chapter}[display]
  {\normalfont\huge\bfseries\color{KURed}}
  {\chaptertitlename\ \thechapter}{20pt}{\Huge}
\titleformat{\section}
  {\normalfont\Large\bfseries\color{KURed!80!black}}
  {\thesection}{1em}{}
\titleformat{\subsection}
  {\normalfont\large\bfseries}
  {\thesubsection}{1em}{}
